\documentclass[main.tex]{subfiles}


\begin{document}

%from pagenumber to up
% \phantomsection 
% \hypertarget{toc}{}



 
   


\section{Специальные записи}

Часто случается, что для решения задачи необходимо сделать \textbf{специальные записи}. Например, задачи на силы решаются на основе рисунков и уравнений сил и~т.\,д. 

Приведем основные такие случаи и укажем соответствующие специальные записи.

\begin{enumerate}
    \item \textbf{Сложное движение.} Если движение тела сложное, то выберите ось, по которой движение тела описать несложно и запишите формулы движения по этой оси.  
    
    \item \textbf{Силы.} Если задача может быть связана с силами, то сделайте рисунок, покажите силы, действующие на тело, и ускорение. Далее запишите уравнения сил по осям, которые видите удобными: $\pm F_1\pm F_2\pm\ldots=\pm ma$.
    
    \item \textbf{Равновесие.} Если задача звязана с покоем <<длинного>> тела (т.\,е. нельзя не учесть его размеры), то сделайте рисунок, покажите силы, действующие на тело. Далее запишите уравнения сил по осям, которые видите удобными: $\pm F_1\pm F_2\pm\ldots=\pm ma$ и запишите уравнение моментов относительно удобной точки: $\pm M_1\pm M_2\pm\ldots=0$.  
    
    \item \textbf{Гидравлический пресс, сообщающиеся сосуды.} Запишите уравнение пресса или сосудов: $P_\text{л}=P_\text{п}$, где $P_\text{л}$ и~$P_\text{п}$~--- давления в левом и правом сосудах.

    \item \textbf{Столкновение (расталкивание) тел.} Если тела сталкиваются и~т.\,д., то сделайте рисунки с телами до и после столкновения, покажите импульсы тел. Далее запишите закон сохранения импульса на соответствующую ось: $\pm p^{\vphantom{'}}_\text{1} \pm p^{\vphantom{'}}_\text{2}\pm\ldots\hm = \pm p'_{1} \pm p'_{2}\pm\ldots$
    
    Если взаимодействие тел упругое, то для системы тел также выполняется закон сохранения полной механической энергии: $E_1=E_2=\ldots$

    \item \textbf{Энергия (механическая).} Если в условии мало данных и/или речь про энергии, то полезно показать рисунки со всеми ключевыми положениями тел. Далее запишите закон сохранения полной механической энергии для тел или системы: $E_1=E_2=\ldots$

    \item \textbf{Смесь газов.} Запишите закон Д\'альтона: $P_\text{см}=P_\text{1\,см}+P_\text{2\,см}+\ldots$

    \item \textbf{Теплообмен.} Если тела <<смешивают>>: одни принимают тепло, другие~--- отдают, то перечислите все процессы с каждым телом (каждому процессу соответствует своя теплота: $Q_1,Q_2,\ldots$). Запишите уравнение теплового баланса: $\pm Q_1\pm Q_2\pm \ldots=0$.

    \item \textbf{Энергия и газ.} Если идет речь о газе в совокупности с энергией/работой, запишите первый закон термодинамики: $Q=\Delta U+A$.


\end{enumerate}




\end{document}
